\documentclass[10pt,a4paper]{amsart}
\usepackage[margin=19mm]{geometry}
\usepackage{amsmath,amssymb,mathtools}
\usepackage[scr=rsfs,cal=euler]{mathalfa}
\usepackage{xfrac}
\usepackage[unicode,hidelinks,bookmarks=false]{hyperref}
\newcommand{\ie}{i.e.\ }
\newcommand{\cf}{cf.\ }
\newcommand{\resp}{respectively\ }
\newcommand{\Subsection}{\subsection}
\providecommand{\nobreakdash}{\nobreak\hbox{-}}
\setlength{\emergencystretch}{2em}
\multlinegap0pt
\usepackage{amssymb}
\usepackage{mathtools}
\newtheorem{lem}{Lemma}
\title{Sharp local estimates for smooth numbers}
\author{Ofir Gorodetsky}
\date{Source: arXiv:2609.10324v1 (CC BY 4.0)}
\begin{document}
\maketitle

\noindent\emph{Source and licence.} Sharp local estimates for smooth numbers. Version arXiv:2609.10324v1 (CC BY 4.0), distributed by arXiv under the Creative Commons Attribution 4.0 International licence, http://creativecommons.org/licenses/by/4.0/, as recorded in the arXiv licence metadata for this exact version and shown on the abstract page https://arxiv.org/abs/2609.10324v1. Full licence notice: \texttt{licenses/arxiv-2609.10324-CC-BY-4.0.txt}.\par\smallskip

\noindent\textit{Editorial scope.} Counted unit: the article's lem lem:alphadiff (source0.tex lines 304--308). The article's complete original proof of it is delivered with it (source0.tex lines 309--315, one \texttt{proof} environment of 7 lines). No mathematics is written, repaired or re-derived: the statement and the proof are the article's own lines, in the article's own notation, and no retained line is substituted or re-flowed.  Retained uncounted as the source-local context the proof needs: lines 290--294.  Nothing was omitted from any retained span, so the package carries no omission record.  No retained line contains a citation, so the wrapper carries no bibliography and no bibliography entry.  No retained line contains a figure environment, an image-inclusion command or drawing code, so no image is delivered and the package declares no asset dependency.  Editorial formatting only, in addition: an AMS \texttt{amsart} preamble that restates the article's own declarations and macros used by the retained lines, namely the environments \texttt{lem} on separate counters, together with the standard packages those lines need.  The retained environments print in this document as Lemma 1 in order of appearance, whereas the article prints the same items with the numbering of its own full document.  The complete native source of the article as distributed in the arXiv e-print for this exact version is bundled unchanged as \texttt{sources/209897/source0.tex}.
\begin{lem}\label{lem:phi}
	Suppose $\log^3 x \ge y \ge 3 \log x$ and $x \ge C$. For all $t$ between $\alpha$ and $\alpha_{\mu^2}$, we have 
	\[\phi^{(2)}(t,y) \asymp \phi^{(2)}_{\mu^2}(t,y) \asymp \log x \log y,\]
	where the derivative is taken with respect to the first variable. If furthermore $y/\log x \to \infty$, then we may replace the  ``$\asymp$'' signs with ``$\sim$''.
\end{lem}

\begin{lem}\label{lem:alphadiff}
	Suppose $\log^3 x \ge y \ge 3\log x$ and $x \ge C$. Then 
	\[\alpha-\alpha_{\mu^2} \asymp \frac{2}{\log x \log y} \sum_{p \le y} \frac{\log p}{p^{2\alpha_{\mu^2}}-1}.\]
	If furthermore $y/\log x \to \infty$, then we may replace the  ``$\asymp$'' sign with ``$\sim$''.
\end{lem}

\begin{proof}
	Comparing the relations $\sum_{p \le y} \log p/(p^{\alpha}-1)=\log x$ and $\sum_{p \le y} \log p/(p^{\alpha_{\mu^2}}+1)=\log x$, we find
	\[\sum_{p \le y} \frac{\log p}{p^{\alpha_{\mu^2}}-1}- \sum_{p \le y} \frac{\log p}{p^{\alpha}-1}  = 2 \sum_{p \le y} \frac{\log p}{p^{2\alpha_{\mu^2}}-1}.\]
	By the mean value theorem, there is a real number $t_{x,y}$ between $\alpha$ and $\alpha_{\mu^2}$ such that
	\[ (\alpha-\alpha_{\mu^2})\phi^{(2)}(t_{x,y},y)  =2 \sum_{p \le y} \frac{\log p}{p^{2\alpha_{\mu^2}}-1}.\]
	We conclude by Lemma \ref{lem:phi}.
\end{proof}

\end{document}
